\section{Critical spin correlations}\label{sec:correlations}

The sector limit of Section~\ref{sec:sectors} and the gluing estimates
of Section~\ref{sec:kernels} now determine the spin correlations after
normalization by the center magnetization.  The microscopic part of
each unit source will be represented by an identical nest of annuli.
Its partition ratio need not have an asymptotic expansion: the same
ratio occurs in the center observable and cancels.

Throughout this section, \(k\geq1\) is fixed,
\(\sigma_i\in\{-1,1\}\), and
\[
 x_{i,n}\in D_n^\circ,\qquad x_{i,n}\longrightarrow z_i\in D,
 \qquad z_i\ne z_j\quad(i\ne j).
\]
Write the continuum Dirichlet Green kernel as
\[
 G_D(z,w)=\frac1{2\pi}\log\frac1{|z-w|}+R_D(z,w).
\]
The function \(R_D\) is smooth across the diagonal on compact subsets
of \(D\times D\).  Define
\[
 H_D(\mathbf z,\boldsymbol\sigma)
 =\sum_{i=1}^k R_D(z_i,z_i)
       +2\sum_{i<j}\sigma_i\sigma_jG_D(z_i,z_j).
\]
No assumption on \(\sum_i\sigma_i\) is imposed.

\subsection{The energy outside small holes}

\begin{lemma}[Exterior energy]\label{correlations:exterior-energy}
Remove the closed disks of radius \(\rho\) about the distinct points
\(z_1,\ldots,z_k\) from \(D\), where \(\rho>0\) is small enough that
these disks are disjoint and contained in \(D\).  For affine gradients
with periods \(2\pi\sigma_i\) about the holes, the minimum energy of
Equation~\eqref{eq:source-3} satisfies
\begin{equation}
 \mathcal M
 =2\pi k\log(1/\rho)
      +4\pi^2H_D(\mathbf z,\boldsymbol\sigma)+O(\rho).
 \tag{12}\label{eq:source-12}
\end{equation}
The error constant may depend on the separated interior configuration,
but is uniform for all sufficiently small \(\rho\).
\end{lemma}

\begin{proof}
Set \(U(y)=\sum_i\sigma_iG_D(y,z_i)\), and rotate
\(-2\pi\nabla U\) by a right angle, with the orientation chosen to give
the indicated periods.  Denote the resulting vector field by \(v\).
It is a smooth affine gradient on the punctured square, is divergence
free there, and is tangent to the outer boundary because \(U\) has
zero Dirichlet trace.  Thus \(v\) is an admissible trial field.

Near \(z_i\), write
\[
 U(y)=\frac{\sigma_i}{2\pi}\log\frac1{|y-z_i|}+u_i(y),
 \qquad
 u_i(z_i)=\sigma_iR_D(z_i,z_i)
                 +\sum_{j\ne i}\sigma_jG_D(z_i,z_j).
\]
The functions \(u_i\) are smooth and harmonic on fixed disjoint
neighborhoods of the sources.  On the boundary of a hole, the radial
singularity contributes no normal component to \(v\).  Its remaining
normal component is bounded uniformly in \(\rho\) and has integral
zero, since it is a tangential derivative of \(u_i\).

Let \(g\in H^1\) on the punctured square.  Integration by parts,
followed by subtraction of the mean of \(g\) on each inner circle,
gives
\[
 \left|\int v\cdot\nabla g\right|
 \leq C\rho\,\|\nabla g\|_2.
\]
Indeed, the \(L^2\) norm of \(v\cdot\nu\) on such a circle is
\(O(\rho^{1/2})\), whereas the trace and Poincar\'e inequalities on
the adjacent annulus of radii \(\rho\) and \(2\rho\) give
\[
 \|g-\text{its circle mean}\|_{L^2(\partial B(z_i,\rho))}
 \leq C\rho^{1/2}
        \|\nabla g\|_{L^2(B(z_i,2\rho)\setminus B(z_i,\rho))}.
\]
The collars are disjoint.  Minimizing
\(\|v+\nabla g\|_2^2\) therefore changes \(\|v\|_2^2\)
by at most \(O(\rho^2)\).

It remains to calculate
\(\|v\|_2^2=(2\pi)^2\int|\nabla U|^2\).
The outer-boundary contribution to Green's identity vanishes.
On the circle about \(z_i\), with \(\nu\) pointing into the hole,
the integral of \(U\partial_\nu U\) is
\[
 \frac1{2\pi}\log(1/\rho)+\sigma_i u_i(z_i)+O(\rho).
\]
The integral of the normal derivative of the regular harmonic part
is zero; in particular it produces no logarithmically enhanced
error.  Summing over the circles gives
Equation~\eqref{eq:source-12}.  These integrations may also be
performed with weak Dirichlet data on the outer square.  Alternatively,
reflection across its sides and corners, away from the sources,
justifies the boundary calculation directly.
\end{proof}

We shall apply this calculation to the rectilinear holes required by
the lattice decomposition.  Let \(P\) be the polygon chosen in
Section~\ref{sec:annuli}, so that
\[
 B(0,1-\varepsilon)\subset P\subset B(0,1+\varepsilon).
\]
Here and in the next lemma, radii are in the scaled coordinates of
\(D\), rather than in lattice units.

\begin{lemma}[Moving lattice holes]\label{correlations:exterior-lattice}
Suppose that \(\rho_n\to\rho>0\), where \(\rho\) is sufficiently
small.  Around \(x_{i,n}\), cut out the interiors in the array
division with hole shape \(\rho_nP\), using the rounding conventions
of Section~\ref{sec:sectors}.  Let \(F_{\rm ext}^s\) be the sector
ratio in the remaining dual graph, with periods
\(2\pi\sigma_i\) about these holes.  Then
\begin{equation}
 \limsup_n
 \left|
 \log F_{\rm ext}^s+
 \frac{2\pi k\log(1/\rho)+4\pi^2H_D(\mathbf z,\boldsymbol\sigma)}
      {2a}
 \right|
 \leq C(\varepsilon+\rho).
 \tag{13}\label{eq:source-13}
\end{equation}
The constant \(C\) depends on the separated limiting configuration
but not on sufficiently small \(\rho\) or \(\varepsilon\).
The same assertion holds after bounded-step mesh rounding.
\end{lemma}

\begin{proof}
Bracket the moving holes by fixed, positively separated rectilinear
holes with dyadic coordinates.  To obtain strict eventual inclusions,
start with circles of radii \((1-3\varepsilon)\rho\) and
\((1+3\varepsilon)\rho\) about the limiting points, and approximate
them and their centers by still smaller dyadic perturbations.
The perturbations can be chosen after fixing \(\rho\) and
\(\varepsilon\).  For all sufficiently large \(n\), the rounded
moving holes lie between these fixed approximations and contain
their respective sources.  The outer boundary remains the same square.

Use one connection before making any of these cuts.  Its restrictions
to the three exteriors have the same winding data.  The cut comparison
in Equation~\eqref{eq:source-2} therefore brackets their sector ratios.
Theorem~\ref{sectors:positive-radius} applies to each fixed polygonal
exterior.  Its continuum minimum energy is in turn bracketed by the
energies for the corresponding circular holes: restriction of an
admissible affine gradient preserves all the remaining periods.
Changing each radius by a factor \(1+O(\varepsilon)\) changes the
logarithmic term of Equation~\eqref{eq:source-12} by
\(O(\varepsilon)\).  That equation now proves the claim.

Every domain used in this argument has positive passages and
separated hole boundaries.  Coherent rounding consequently preserves
connectivity and the cycles defining the periods.  In particular the
argument uses only fixed-domain limits; it does not require uniform
convergence under arbitrary motion of the lattice marks.
\end{proof}

\subsection{Identical nests and their exterior attachment}

Fix the annular parameters \(\varepsilon,\lambda,R_0\) as in
Section~\ref{sec:kernels}, and write
\[
 R_j=R_0\lambda^j,\qquad
 N=N(n)=\max\{j:R_j\leq\eta n\},
\]
where \(\eta>0\) is small.  The \(R_j\) are now lattice radii.
Place identical translated nests, from their filled cores through
their \(N\)-th annuli, around all \(x_{i,n}\).
For sufficiently small \(\eta\), these nests have disjoint closures,
are contained in the spin square, and leave a connected exterior.
Use the same parameters and the same \(N\) for a single nest centered
at \(0\), which will represent the observable \(a_n\).
Since \(nx_{i,n}\in\Z^2\), the microscopic cores and all annuli are
exact lattice translates.  The restricted connections differ only
by integral gauge transformations; reversing a unit period also
leaves a bare nest ratio unchanged.

We may pass further along any subsequence so that
\[
 R_N/n\longrightarrow t\in[\eta/\lambda,\eta].
\]
For fixed \(\eta,\lambda\), this is a positive macroscopic radius.
The following estimate concerns only the joint kernel which attaches
the last annuli to the exterior.  Its notation is that of
Lemma~\ref{kernels:bounds}: \(d_{\rm ext}\) is the centered kernel
at zero relative guard data, and \(m_{N,i}^s\) is the cut law of
the last guard in the \(i\)-th annulus.

\begin{lemma}[Exterior attachment]\label{correlations:attachment}
For every sufficiently small \(\xi>0\), the parameters may be chosen
so that Equation~\eqref{eq:source-11} holds throughout the regular
annuli and
\begin{equation}
 \limsup_n\int
       \left(1-\frac{T_{\rm ext}^s}{d_{\rm ext}}\right)
                 \prod_{i=1}^k dm_{N,i}^s
 \leq\xi .
 \tag{14}\label{eq:source-14}
\end{equation}
This holds both for the shifted and centered calculations and for
the single center nest.  The parameters depend on the fixed
separated configuration and on \(\xi\), but not on \(n\).
The errors in Equation~\eqref{eq:source-13} at the radii in this
construction may simultaneously be made at most \(\xi\).
\end{lemma}

\begin{proof}
For the centered calculation, Equation~\eqref{eq:source-9} bounds
the attachment loss by the sum of the flat-pin interactions in
the incident last annuli.  These interactions can be made as small
as desired by the choices in Section~\ref{sec:annuli}.
The bound depends only on these annuli, not on the exterior.

For the shifted calculation use the exact identity
in Equation~\eqref{eq:source-10}, now for the last annuli together
with their common exterior.  Restoring them produces an exterior
whose smaller holes have radius parameter tending to \(t/\lambda\).
Before restoration, the factors are the \(k\) last-annulus sector
ratios and the exterior ratio with holes of radius parameter \(t\).
By Theorem~\ref{sectors:positive-radius}, each annular energy is
\(2\pi\log\lambda+O(\varepsilon)\).  By
Equation~\eqref{eq:source-13}, the restored and cut exterior energies
have the same term \(4\pi^2H_D(\mathbf z,\boldsymbol\sigma)\).
The difference between their radius terms is exactly
\(k\,2\pi\log\lambda\).  Consequently the logarithm of the
restored sector ratio divided by the product of the cut ratios
has absolute limiting error at most \(C(\varepsilon+\eta)\).
The circular energy brackets make this error uniform over
\(t\in[\eta/\lambda,\eta]\) and over sufficiently large fixed
\(\lambda\).

Choose \(\varepsilon\) and \(\eta\) small, then choose \(\lambda\)
large enough for the flat-pin interactions, and finally choose
\(R_0\) large enough for the fixed annular geometry.  These choices
retain the regular-interface estimate
Equation~\eqref{eq:source-11}.  The preceding centered estimate,
the restored-to-cut identity, and the upper bound
\(T_{\rm ext}^s/d_{\rm ext}\leq1\) now give
Equation~\eqref{eq:source-14}, after making the preliminary errors
small enough in terms of \(\xi\).  Apply the same choices to the
one-source configuration at \(0\).

The argument was given after extraction of a limit \(t\).
Every subsequence has such a further subsequence, and the error
bounds are uniform over its possible values.  It therefore proves
the stated limiting bounds on the full sequence.  All height and
pin limits used here precede variation of the fixed geometric
parameters; no rate uniform in a shrinking physical radius has
been assumed.
\end{proof}

\begin{proposition}[Center-normalized correlations]
\label{correlations:separated}
For the separated moving marks fixed at the beginning of this section,
\begin{equation}
 \lim_n\log\left(
 \frac{\E_n^0\exp(i\sum_i\sigma_i\theta_{x_{i,n}})}{a_n^k}
 \right)
 =
 -\frac{H_D(\mathbf z,\boldsymbol\sigma)-kR_D(0,0)}
        {2K_*}.
 \tag{15}\label{eq:source-15}
\end{equation}
\end{proposition}

\begin{proof}
Let \(F_{\rm nest}(R_N)\) denote the bare sector ratio of one whole
free nest with a unit source, including its core.  It is the same
for all the translated sources, for either sign, and for the
center reference.  By Lemma~\ref{kernels:multiplication}, the
partition function of the whole graph is the product of its cut
block partition functions times its gluing-kernel expectation.
After the kernels internal to each nest have been integrated, its
last-guard law has density at most \(2\) relative to \(m_N^s\),
for all sufficiently large \(N\), by
Corollary~\ref{kernels:last-guard}.

Write \(e_s\) for the expectation of
\(T_{\rm ext}^s/d_{\rm ext}\) under the product of these propagated
last-guard laws, and \(e_0\) for the corresponding centered
expectation.  Equation~\eqref{eq:source-14} and the product density
bound give
\[
 1-2^k\xi-o(1)\leq e_s\leq1
\]
in both sectors.  The centered factor \(d_{\rm ext}\) cancels
between the numerator and denominator, leaving the exact identity
\[
 \E_n^0\exp\left(i\sum_i\sigma_i\theta_{x_{i,n}}\right)
 =F_{\rm nest}(R_N)^k F_{\rm ext}^s\,\frac{e_s}{e_0}.
\]
Taking logarithms, with \(\xi\) small at this fixed \(k\), gives
\[
 \log\E_n^0\exp\left(i\sum_i\sigma_i\theta_{x_{i,n}}\right)
 =k\log F_{\rm nest}(R_N)+\log F_{\rm ext}^s
                         +O_k(\xi)+o(1).
\]
The same calculation with the single center source expresses
\(\log a_n\) using this identical nest ratio.  Subtracting \(k\)
times that expression cancels the microscopic term exactly.
In Equation~\eqref{eq:source-13}, the common-radius logarithms
also cancel.  The remaining coefficient is
\[
 \frac{4\pi^2}{2a}=\frac1{2K_*},
 \qquad a=8\pi,\quad K_*=\frac2\pi.
\]
This proves Equation~\eqref{eq:source-15} up to an error that can
be made arbitrarily small.

For completeness, the choices of nests may depend on that accuracy.
This is harmless: the quantity on the left is always the same exact
observable.  If its asserted convergence failed, a subsequence with
a fixed positive error would, at a sufficiently high chosen accuracy,
have a further subsequence with \(R_N/n\to t\) and a smaller error
by the preceding calculation, a contradiction.
\end{proof}

\subsection{The deterministic diagonal correction}

\begin{lemma}[Green diagonal differences]
\label{correlations:green-diagonal}
If \(x_n\in D_n^\circ\) and \(x_n\to z\in D\), then
\begin{equation}
 G_n(x_n,x_n)-G_n(0,0)
       \longrightarrow R_D(z,z)-R_D(0,0).
 \tag{16}\label{eq:source-16}
\end{equation}
These differences are bounded uniformly when \(x_n\) ranges over
a fixed compact subset of \(D\).
\end{lemma}

\begin{proof}
The killed Green estimate in \cite[Lemma~4.2]{B} gives local uniform
convergence of \(G_n\) to \(G_D\) off the diagonal, with exactly the
unscaled Laplacian convention used here.  We explain how to recover
the difference of the diagonals without specifying their common
lattice constant.

Choose a small fixed box \(Q\) of displacement vectors about \(0\)
such that \(x_n+Q\) and \(Q\) remain in \(D\) for all large \(n\).
For lattice displacement vectors \(v\), set
\[
 h_n(v)=G_n(x_n,x_n+v)-G_n(0,v).
\]
Both kernels have the same unit source at \(v=0\).
Their difference is therefore discrete harmonic throughout \(Q\),
including at \(v=0\).  On a coherently rounded boundary of \(Q\),
the off-diagonal convergence identifies its limit as
\[
 G_D(z,z+v)-G_D(0,v)
       =R_D(z,z+v)-R_D(0,v).
\]
The discrete maximum principle bounds \(h_n(0)\) between the
minimum and maximum of these boundary values, up to an error
tending to zero.  Shrinking \(Q\) and using continuity gives
Equation~\eqref{eq:source-16}.
If uniform boundedness on an interior compact failed, a violating
sequence would have a convergent subsequence of its locations.
The convergence just proved would contradict that violation.
The finitely many remaining lattice sizes cause no difficulty.
\end{proof}

For \(\sigma\in\{-1,1\}\), define the normalized point fields
\[
 W_n^\sigma(x)
 =a_n^{-1}
   \exp\left\{\frac{G_n(x,x)-G_n(0,0)}{2K_*}\right\}
       e^{i\sigma\theta_x}.
\]
Thus \(S_n(f)=n^{-2}\sum_x f(x)W_n^+(x)\), and changing
the sign conjugates the point field without conjugating a test
function.

\begin{corollary}[Normalized separated correlations]
\label{correlations:normalized}
For every fixed separated moving configuration as above,
\begin{equation}
 \lim_n\E_n^0\prod_{i=1}^kW_n^{\sigma_i}(x_{i,n})
 =\exp\left\{-\frac1{K_*}
          \sum_{i<j}\sigma_i\sigma_jG_D(z_i,z_j)\right\}.
 \tag{17}\label{eq:source-17}
\end{equation}
The convergence to the same kernel evaluated at the lattice
locations is uniform when those locations range over a compact
interior set and their mutual distances are bounded below by a
fixed positive number.
\end{corollary}

\begin{proof}
Exponentiate Equation~\eqref{eq:source-15}, multiply by the
deterministic normalizing factors, and use
Equation~\eqref{eq:source-16}.
The terms
\(\sum_iR_D(z_i,z_i)-kR_D(0,0)\) cancel exactly.
For uniformity, a sequence of violations in a compact separated
configuration space would have a convergent subsequence, to which
the just-proved sequential convergence applies.
\end{proof}

\subsection{A bound through the two-point diagonal}

The separated formula alone does not control moment sums near
collisions.  For that purpose a weaker, integrable power bound
suffices.  We obtain it by cutting two disjoint nests at a scale
comparable to the distance between their marks.

Fix the annular parameters once and for all, with accuracy high
enough that the center attachment in
Lemma~\ref{correlations:attachment} is bounded away from zero.
Keep \(N(n)=\max\{j:R_j\leq\eta n\}\) for these fixed parameters.

\begin{lemma}[Nest lower bounds]\label{correlations:nest-bounds}
There are constants \(c,c'>0\), independent of \(n,N,j\), such that,
for all sufficiently large \(n\),
\begin{equation}
 a_n\geq c\,F_{\rm nest}(R_{N(n)}),
 \tag{18}\label{eq:source-18}
\end{equation}
and, for every \(N\geq j\geq0\),
\begin{equation}
 F_{\rm nest}(R_N)
       \geq c'\lambda^{-(N-j)/2}F_{\rm nest}(R_j).
 \tag{19}\label{eq:source-19}
\end{equation}
Here \(j=0\) means the filled core alone.
\end{lemma}

\begin{proof}
For the single center source, the normalized exterior attachment
has a positive lower bound in the shifted sector and is at most
one in the centered sector.  Its exterior sector ratio is also
bounded away from zero by Equation~\eqref{eq:source-13}:
every subsequential hole radius lies in the fixed interval
\([\eta/\lambda,\eta]\).  The exact gluing identity therefore
gives Equation~\eqref{eq:source-18}.  The compact subsequence
argument makes the lower bound uniform for all sufficiently
large \(n\).

Next add one regular annulus to a nest.  Its sector ratio
multiplies the nest ratio, together with the ratio of the
shifted and centered attachment expectations.  Normalize these
two expectations by their common centered zero-data value.
The centered expectation is at most one.  Once the nest is
sufficiently long, its last-data density is at most two, so
Equation~\eqref{eq:source-11} makes the shifted expectation at
least \(1-2\xi\).

By Theorem~\ref{sectors:positive-radius} and the circular annular
energy bracket, the annular sector factor at large lattice scales
is bounded below by
\[
 \lambda^{-1/8}\exp\{-O(\varepsilon)-o(1)\},
 \qquad \frac{\pi}{a}=\frac18.
\]
Include in the fixed parameter choices that \(\lambda\) and then
the lattice threshold are sufficiently large.  Each sufficiently
late extension factor is then at least \(\lambda^{-1/2}\).
There are only finitely many earlier extension factors, all
strictly positive.  Their losses relative to this lower bound
can be absorbed into a single \(c'>0\), valid for every interval
of extension indices.  Multiplication proves
Equation~\eqref{eq:source-19}.
\end{proof}

\begin{proposition}[Two-point collision bound]
\label{correlations:diagonal}
For every compact \(E\Subset D\), there is \(C_E<\infty\) such that,
for all \(n\geq1\), all \(x,y\in E\cap D_n^\circ\), including \(x=y\),
and all \(\sigma,\tau\in\{-1,1\}\),
\begin{equation}
 0\leq\E_n^0\bigl[W_n^\sigma(x)W_n^\tau(y)\bigr]
       \leq C_E\bigl(|x-y|+1/n\bigr)^{-1}.
 \tag{20}\label{eq:source-20}
\end{equation}
Expectations of products of any number of the point fields
\(W_n^\sigma\), with arbitrary signs and repeated locations, are
also nonnegative.
\end{proposition}

\begin{proof}
Write \(r=n|x-y|\).  Choose \(c_E>0\) so small that nests of
radius \(R\leq c_Er\), centered at \(x\) and \(y\) in lattice
units, are disjoint and contained in the spin square.  Decrease
\(c_E\), if necessary, so that \(c_Er\leq\eta n\) for all
locations under consideration.

If there is a largest \(j\geq0\) with \(R_j\leq c_Er\), cut
out these two nests.  The restriction of the magnetic sector
to each nest has its single unit source.  The cut comparison
and the upper bound one for the remaining sector ratio give
\[
 \E_n^0e^{i\sigma\theta_x+i\tau\theta_y}
       \leq F_{\rm nest}(R_j)^2.
\]
Equations~\eqref{eq:source-18} and \eqref{eq:source-19} yield
\[
 \frac{\E_n^0e^{i\sigma\theta_x+i\tau\theta_y}}{a_n^2}
 \leq C\lambda^{N-j}.
\]
The maximality of \(j\) gives \(R_j>c_Er/\lambda\), while
\(R_N\leq\eta n\).  Since \(R_N/R_j=\lambda^{N-j}\),
this is at most \(C_E n/(1+r)\), after increasing the constant.

If no such \(j\) exists, then \(r<R_0/c_E\).  Use the bare
expectation bound by one and Equation~\eqref{eq:source-19}
with \(j=0\), together with Equation~\eqref{eq:source-18}.
The fixed core ratio is positive, and hence
\[
 a_n^{-2}\leq C\lambda^N\leq Cn
       \leq C_E\frac{n}{1+r}.
\]
This case includes repeated locations with either pair of signs.
The deterministic Green factors in the two \(W_n\)'s are uniformly
bounded on \(E\) by Lemma~\ref{correlations:green-diagonal}.
They therefore preserve the bound.  Increasing \(C_E\) handles
the finitely many smaller lattice sizes.

Finally, the Fourier current expansion of the pinned XY law
expresses each integer-charge correlation as a sum of
nonnegative weights divided by a positive partition function.
The fixed boundary imposes no neutrality constraint; repeated
locations merely add their charges.  Multiplication by the
positive deterministic factors proves all the asserted
nonnegativity statements.
\end{proof}
